
\chapter*{Conclusiones y recomendaciones}
\addcontentsline{toc}{chapter}{Conclusiones y recomendaciones}

En este trabajo se estudió el vacío de la intersección de una gramática
de concatenación de rango con un lenguaje regular finito como marco para
el problema de la satisfacibilidad booleana.

En el capítulo~\ref{cap:rcg-sat} se presentaron las gramáticas de
concatenación de rango y los antecedentes en que el SAT se aborda como
la intersección de un lenguaje regular de interpretaciones con una
gramática que impone la consistencia entre las ocurrencias de cada
variable, en SATCF y SATSRC.

En el capítulo~\ref{cap:vacio-interseccion-rcg-regular} se definió el
problema del vacío de la intersección de una RCG con un lenguaje regular
finito y se dio un algoritmo de decisión, en el que se enumeran las
cadenas del autómata y se reconoce cada una con la gramática. Con él se
decide el problema, pero su costo es exponencial cuando la cantidad de
cadenas aceptadas por el autómata es exponencial, y polinomial cuando esa
cantidad está acotada por un polinomio.

En el capítulo~\ref{cap:modelacion-sat} ese resultado se aplicó al SAT: a
cada fórmula booleana se le asocian un autómata booleano y una gramática
de consistencia, de modo que la fórmula es satisfacible si y solo si la
intersección de sus lenguajes no es vacía. Decidir ese vacío es el SAT,
y para una RCG con un autómata finito acíclico ese problema es
NP-completo~\cite{bertsch2001rcg}.

En el capítulo~\ref{cap:interseccion} se construyó la gramática producto
$G'$ de una RCG $G$ y un autómata finito acíclico $A$, con
$L(G') = L(G) \cap L(A)$, y se demostró su correctitud. La gramática
producto es una RCG sobre la que estudiar el vacío de la intersección,
pero su tamaño puede ser exponencial: para la gramática de consistencia y
el autómata booleano de una fórmula con $n$ ocurrencias de variable
booleana, $|G'| = 2^{O(n\log n)}$, y las cláusulas de la consistencia son
el término dominante.

En conjunto, estos resultados son un marco uniforme para el SAT en torno
a una sola noción: el vacío de la intersección de una RCG con un lenguaje
regular finito.

La gramática de consistencia del capítulo~\ref{cap:modelacion-sat} tiene
aridad uno, para cualquier fórmula, gracias a la no linealidad. Para
abarcar más fórmulas que SATSRC sin que
la aridad crezca con los cruces entre ocurrencias hace falta la no
linealidad, pues sin ella se reconocen los mismos lenguajes que con una
sRCG, como se vio en el capítulo~\ref{cap:rcg-sat}. Por esa misma no
linealidad la modelación queda fuera de las subclases conocidas con vacío
de la intersección polinomial, como SATCF y SATSRC, y se pierde la
garantía directa de que el vacío sea polinomial. Si solo se modelara una subclase
de fórmulas, el vacío podría ser polinomial sin que ello implique
$P = NP$; decidirlo en tiempo polinomial para todas las fórmulas sí
implicaría $P = NP$.

La gramática producto del capítulo~\ref{cap:interseccion} es el objeto
sobre el que abordar ese vacío: hallar un método polinomial para todas
las fórmulas, probar que no existe, o hallarlo solo al restringir la
modelación a una subclase. Para decidir el vacío en tiempo polinomial
sobre la gramática producto hace falta eliminar antes sus predicados y
cláusulas no productivos.

A partir del trabajo realizado se proponen las siguientes líneas de
investigación:

\begin{itemize}
  \item Eliminar de la gramática producto los predicados y cláusulas no
    productivos y analizar la complejidad de decidir el vacío de la
    intersección sobre la gramática resultante.
  \item Buscar subclases no lineales de RCG cuyo vacío de la
    intersección con un lenguaje regular finito se decida en tiempo
    polinomial, y estudiar si la gramática de consistencia se adapta a
    alguna de ellas.
  \item Modificar las RCG para obtener un formalismo cuyo vacío de la
    intersección con un lenguaje regular finito se decida en tiempo
    polinomial y con el que se exprese la consistencia de fórmulas
    distintas de las que ya se modelan con SATCF y SATSRC.
  \item Caracterizar las fórmulas booleanas cuyo autómata booleano tiene
    una cantidad de caminos aceptantes acotada por un polinomio, donde el
    SAT se decide en tiempo polinomial con la enumeración del
    capítulo~\ref{cap:vacio-interseccion-rcg-regular}.
\end{itemize}

